% Frozen source-aware proof packet, 2026-08-25.
% This file has not been submitted to an external verifier.
%
% Authenticated primary-source dependencies:
% (1) Leonard Pille-Schneider, Global pluripotential theory on hybrid
%     spaces, arXiv:2209.04879v2, local PDF SHA-256
%     55defd54d8201c94d08f727e76f551c759275c68cde354570cb41e7a02df54f7.
%     Used: Proposition 2.25, Definition 3.9, Theorems 3.10 and 3.13,
%     Proposition 3.16, Lemma 3.17, and Question 3.18.
% (2) Charles Favre, Degeneration of endomorphisms of the complex
%     projective space in the hybrid space, arXiv:1611.08490v2,
%     local PDF SHA-256
%     137ddb0245b4f09ff8b7a4a2880bd3846f146cd013f2ce9563d7f90f45ea1a5f.
%     Used: Definition 4.1 and its following remark (printed pp. 25--26),
%     and Question 1 (printed p. 33).
%
% Scope firewalls:
% * This is the canonical psh gluing of Pille-Schneider's Theorem 3.13,
%   not arbitrary set-theoretic gluing.
% * Favre's literal Question 1 is checked independently; no equivalence
%   with Pille-Schneider's Question 3.18 is assumed.
% * The example is a base-scalar obstruction.  It does not decide a
%   strengthened question with a normalization controlling base scalars.

\begin{theorem}
Fix a hybrid radius \(r_0\in(0,1)\).  There exist a smooth trivial
polarized degeneration
\[
 \pi\colon X=\mathbf P^1_{\mathbf C}\times\mathbb D^*
       \longrightarrow\mathbb D^*,
 \qquad
 L=\operatorname{pr}_1^*\mathcal O_{\mathbf P^1}(1),
\]
and a smooth semipositive metric \(\phi\) on \(L\), with logarithmic
growth at the origin, having the following properties.
\begin{enumerate}
\item The metric \(\phi\) extends as an everywhere finite psh metric on
the product model
\[
 \mathcal X=\mathbf P^1_{\mathbf C}\times\mathbb D,
 \qquad
 \mathcal L=\operatorname{pr}_1^*\mathcal O_{\mathbf P^1}(1).
\]
\item Its associated non-Archimedean metric is the continuous product
Fubini--Study model metric:
\[
                   \phi^{\mathrm{NA}}
                   =\phi_{\mathrm{FS}}^{\mathrm{NA}}.
\]
\item The canonical hybrid psh metric \(\phi^{\mathrm{hyb}}\), obtained
by putting \(\phi\) on the complex fibres and
\(\phi^{\mathrm{NA}}\) on the non-Archimedean fibre, is not continuous.
More precisely, for every \(p\in\mathbf P^1(\mathbf C)\) there are two
sequences \(x_m^+\) and \(x_m^-\) in the complex locus that both
converge to the same point \(p_{\mathbf C((t))}\) and for which
\[
\begin{aligned}
 \lim_{m\to\infty}
 \bigl(\phi^{\mathrm{hyb}}
       -\phi_{\mathrm{FS}}^{\mathrm{hyb}}\bigr)(x_m^+)
   &=\log r_0,\\
 \lim_{m\to\infty}
 \bigl(\phi^{\mathrm{hyb}}
       -\phi_{\mathrm{FS}}^{\mathrm{hyb}}\bigr)(x_m^-)
   &=0.
\end{aligned}
\]
\end{enumerate}
Consequently, continuity of a semipositive metric on the punctured total
space together with continuity of its associated non-Archimedean metric
does not imply continuity of the canonical hybrid psh extension.

Moreover, relative to the product Fubini--Study metric, the resulting
potential satisfies separately all three conditions in Favre's
Question~1: it is continuous near every punctured
\(\overline X_r\), its curvature is the restriction of a positive
closed \((1,1)\)-current on a model neighborhood, and its
generic-Lelong function extends continuously to the non-Archimedean
analytification.  Nevertheless, that potential is not uniform in the
sense of Favre's Definition~4.1.
\end{theorem}

\begin{proof}
For \(j\geq1\), put
\[
 s_j=2^j,\qquad
 \rho_j=\exp(-s_j),\qquad
 c_j=4^{-j},\qquad
 \epsilon_j=\exp(-s_j/c_j)=\exp(-8^j),
\]
and define on the unit disk
\[
 q_j(t)=\frac{c_j}{2}
   \log\frac{|t-\rho_j|^2+\epsilon_j^2}{5}.
\tag{1}\label{eq:p18-qj}
\]
The functions \(q_j\) are smooth and subharmonic.  Indeed,
\(\log(|t-a|^2+\epsilon^2)\) has nonnegative
\(\mathrm d\mathrm d^c\), and multiplication by \(c_j/2>0\)
preserves this property.  Moreover \(q_j\leq0\) on \(\mathbb D\):
for \(|t|<1\), the numerator in \eqref{eq:p18-qj} is strictly smaller
than \(5\).

Set
\[
                         u(t)=\sum_{j=1}^{\infty}q_j(t).
\tag{2}\label{eq:p18-u}
\]
Let \(K\Subset\mathbb D^*\).  For all sufficiently large \(j\), the
points \(\rho_j\) and the disks of radius \(\epsilon_j\) about them
have a fixed positive distance from \(K\).  Every derivative, on \(K\),
of the logarithm occurring in \eqref{eq:p18-qj} is then bounded by a
constant depending only on \(K\) and the order of the derivative.
Since \(\sum_jc_j<\infty\), the series in \eqref{eq:p18-u} converges in
\(C^\infty(K)\).  Hence \(u\) is smooth and subharmonic on
\(\mathbb D^*\).

At the origin one has
\[
\begin{aligned}
 q_j(0)
 &=
 \frac{c_j}{2}
 \left(
  2\log\rho_j-\log5+
  \log\left(1+\frac{\epsilon_j^2}{\rho_j^2}\right)
 \right)\\
 &=-c_js_j-\frac{c_j}{2}\log5+O(c_j)
  =-2^{-j}+O(4^{-j}).
\end{aligned}
\tag{3}\label{eq:p18-at-zero}
\]
Thus \(\sum_jq_j(0)>-\infty\).  The partial sums in
\eqref{eq:p18-u} form a decreasing sequence of subharmonic functions on
\(\mathbb D\), because every \(q_j\leq0\).  Their limit is not
identically \(-\infty\), by \eqref{eq:p18-at-zero}; therefore the
decreasing-limit theorem gives a subharmonic extension of \(u\) to
\(\mathbb D\), with the finite value prescribed by
\eqref{eq:p18-at-zero}.

This extension has zero Lelong number at the origin.  To see this without
using any pointwise radial assertion, let \(M_u(r)\) be its circular
mean.  The submean inequality and \(u\leq0\) give
\[
                         u(0)\leq M_u(r)\leq0.
\]
The circular-mean characterization of the one-variable Lelong number
therefore yields
\[
                       \nu_0(u)
                       =\lim_{r\to0}\frac{M_u(r)}{\log r}=0.
\tag{4}\label{eq:p18-zero-lelong}
\]

We next establish the logarithmic-growth estimate and isolate two
different normalized limits.  Let
\[
                B_j=\{\,t\in\mathbb D:
                         |t-\rho_j|<\rho_j/2\,\}.
\]
The disks \(B_j\) are pairwise disjoint.  Indeed,
\(\rho_{j+1}=\rho_j^2\), \(\rho_1<1/3\), and hence
\[
                    \frac32\rho_{j+1}
                    <\frac12\rho_j .
\]
If \(t\notin B_j\), then
\[
\begin{aligned}
 q_j(t)
 &\geq\frac{c_j}{2}
       \log\frac{\rho_j^2}{20}\\
 &=-c_js_j-\frac{c_j}{2}\log20.
\end{aligned}
\tag{5}\label{eq:p18-outside-well}
\]
If \(t\) belongs to none of the \(B_j\), summing
\eqref{eq:p18-outside-well} gives \(u(t)\geq-C\), since
\[
                \sum_{j\geq1}c_js_j=\sum_{j\geq1}2^{-j}<\infty.
\]
If \(t\in B_m\), pairwise disjointness permits the use of
\eqref{eq:p18-outside-well} for all \(j\ne m\), while
\[
\begin{aligned}
 q_m(t)
 &\geq\frac{c_m}{2}\log\frac{\epsilon_m^2}{5}\\
 &=-s_m-\frac{c_m}{2}\log5.
\end{aligned}
\tag{6}\label{eq:p18-inside-well}
\]
Moreover \(-\log|t|=s_m+O(1)\) uniformly for \(t\in B_m\).
Equations \eqref{eq:p18-outside-well}--\eqref{eq:p18-inside-well},
together with \(u\leq0\), prove that, after increasing \(C\),
\[
             -\log|t|^{-1}-C\leq u(t)\leq0
             \qquad(0<|t|<1).
\tag{7}\label{eq:p18-log-growth}
\]

At \(t=\rho_m\), the \(m\)-th summand is exactly
\[
                  q_m(\rho_m)
                  =-s_m-\frac{c_m}{2}\log5.
\]
Since \(\rho_m\notin B_j\) for \(j\ne m\), the sum of all other
summands lies between \(-C\) and \(0\), by
\eqref{eq:p18-outside-well}.  It follows that
\[
                         u(\rho_m)=-s_m+O(1).
\tag{8}\label{eq:p18-plus-sequence}
\]
On the other hand, \(-\rho_m\) belongs to none of the \(B_j\), because
\(|-\rho_m-\rho_j|=\rho_m+\rho_j>\rho_j/2\).  Hence
\[
                         -C\leq u(-\rho_m)\leq0.
\tag{9}\label{eq:p18-minus-sequence}
\]
Since \(\log|\rho_m|=-s_m\), equations
\eqref{eq:p18-plus-sequence}--\eqref{eq:p18-minus-sequence} give
\[
\begin{aligned}
 \frac{\log r_0}{\log|\rho_m|}\,u(\rho_m)
   &\longrightarrow\log r_0,\\
 \frac{\log r_0}{\log|-\rho_m|}\,u(-\rho_m)
   &\longrightarrow0.
\end{aligned}
\tag{10}\label{eq:p18-two-limits}
\]

We now put this scalar example on a polarized family.  Let
\[
 \mathcal X=\mathbf P^1_{\mathbf C}\times\mathbb D,\qquad
 \mathcal L=\operatorname{pr}_1^*\mathcal O_{\mathbf P^1}(1),
\]
and denote by \(\phi_{\mathrm{FS}}\) the product Fubini--Study model
metric.  On the restriction \(X=\mathcal X|_{\mathbb D^*}\), set
\[
                       \phi=\phi_{\mathrm{FS}}+u\circ\pi .
\tag{11}\label{eq:p18-product-metric}
\]
The metric is smooth on \(X\), and its curvature is
\[
 \mathrm d\mathrm d^c\phi
 =
 \operatorname{pr}_1^*\omega_{\mathrm{FS}}
 +\pi^*\mathrm d\mathrm d^c u
 \geq0.
\tag{12}\label{eq:p18-positive-curvature}
\]
It is therefore a continuous semipositive metric on the punctured total
space.  The subharmonic extension of \(u\) constructed above makes
\eqref{eq:p18-product-metric} an everywhere finite psh metric on
\(\mathcal L\) over \(\mathcal X\).  This proves logarithmic growth by
the first equivalent condition in Pille-Schneider's Definition~3.9;
\eqref{eq:p18-log-growth} also verifies the two-sided logarithmic bound
directly.

We compute the associated non-Archimedean metric.  Let \(E\) be a prime
vertical divisor on an arbitrary model dominating \(\mathcal X\).  At a
general smooth point of \(E\), one may write
\[
                          t=\zeta^{b_E}h,
\]
where \(\zeta=0\) is a local equation for \(E\), \(h\) is a unit, and
\(b_E=\operatorname{ord}_E(t)\).  In one complex variable the Lelong
number is the atomic mass at the origin of the Riesz measure, up to the
fixed normalization.  Pullback by a holomorphic map of order \(b_E\)
multiplies this atomic mass by \(b_E\).  Thus
\[
                 \nu_E(u\circ\pi)
                 =b_E\nu_0(u)=0
\tag{13}\label{eq:p18-all-divisorial-zero}
\]
by \eqref{eq:p18-zero-lelong}.  The smooth Fubini--Study reference has
zero relative generic Lelong numbers as well.  Pille-Schneider's
Theorem~3.10 therefore gives
\[
                  \phi^{\mathrm{NA}}
                  =\phi_{\mathrm{FS}}^{\mathrm{NA}},
\tag{14}\label{eq:p18-na-model}
\]
the continuous product model metric.  Equivalently, the relative
potential is zero on every divisorial point, and the unique extension in
that theorem is the zero model function.

Pille-Schneider's Theorem~3.13 now applies to
\eqref{eq:p18-product-metric}: the metric \(\phi\) and the endpoint
\eqref{eq:p18-na-model} form one canonical hybrid psh metric
\(\phi^{\mathrm{hyb}}\).  This invokes genuine hybrid psh gluing, rather
than an arbitrary assignment of values on two disjoint sets.

Fix \(p\in\mathbf P^1(\mathbf C)\), and write \(p_{\mathbf C((t))}\)
for the point on the non-Archimedean fibre induced by the constant
algebraic section.  Functoriality of hybrid analytification, applied to
that section, gives
\[
\begin{aligned}
 x_m^+&=(p,\rho_m)\longrightarrow p_{\mathbf C((t))},\\
 x_m^-&=(p,-\rho_m)\longrightarrow p_{\mathbf C((t))}.
\end{aligned}
\tag{15}\label{eq:p18-hybrid-convergence}
\]
The hybrid metric induced by the product Fubini--Study model is
continuous.  In additive weights, its difference from
\(\phi^{\mathrm{hyb}}\) on the fibre over \(t\ne0\) is
\[
 \bigl(\phi^{\mathrm{hyb}}
       -\phi_{\mathrm{FS}}^{\mathrm{hyb}}\bigr)|_{X_t}
 =\frac{\log r_0}{\log|t|}\,u(t),
\tag{16}\label{eq:p18-hybrid-normalization}
\]
whereas its difference on the central non-Archimedean fibre is zero by
\eqref{eq:p18-na-model}.  Combining
\eqref{eq:p18-two-limits},
\eqref{eq:p18-hybrid-convergence}, and
\eqref{eq:p18-hybrid-normalization} proves the two limits in the
statement.  Since \(\log r_0\ne0\), the canonical hybrid psh metric is
not continuous.

It remains to check Favre's formulation independently.  Relative to the
product Fubini--Study metric, take
\[
                         \varphi=u\circ\pi
\]
and, on the product model, take the positive closed current
\[
 T=
 \operatorname{pr}_1^*\omega_{\mathrm{FS}}
 +\pi^*\mathrm d\mathrm d^c u.
\tag{17}\label{eq:p18-favre-current}
\]
The potential is smooth on \(X\), hence continuous near every punctured
\(\overline X_r\).  Equation \eqref{eq:p18-positive-curvature} says that
\eqref{eq:p18-favre-current} restricts to
\(\omega_{\mathrm{FS}}+\mathrm d\mathrm d^c\varphi\), and
\eqref{eq:p18-all-divisorial-zero} says that Favre's generic-Lelong
function is \(g_T\equiv0\).  It therefore extends continuously to the
non-Archimedean analytification.  All three bullets of Favre's
Question~1 hold.

Suppose, for contradiction, that \(\varphi\) were uniform.  The remark
following Favre's Definition~4.1 states that its normalized function on
the punctured hybrid locus would then extend continuously to the hybrid
space.  For the hybrid radius appearing in that definition, the two
sequences in \eqref{eq:p18-hybrid-convergence} have the distinct limits
computed in \eqref{eq:p18-two-limits}.  Such a continuous extension is
impossible.  Thus \(\varphi\) is not uniform, completing the proof.
\end{proof}
